\begin{task}
Wyznacz spadek napięcia na rezystorze $R_1$
\begin{center}\begin{circuitikz}[european] 
\ctikzset{voltage/distance from line=0.2} 
\ctikzset{voltage/european label distance=2.00}  
\draw
 (0,0) to[R,l=$R_1$] (2,0) 
       to[R,l=$R_2$] (2,-2)
       to[short] (0,-2)
       to[american current source,l=$E$] (0,0) 
;
\end{circuitikz}
\end{center}
\subsubsection{Rozwiązanie}
Pokazany układ możemy uprościć poprzez szeregowe połączenie rezystorów $R_1$ i $R_2$. Rezystor zastępczy będzie miał rezystancję $R_z = R_1+R_2$.

\begin{center}\begin{circuitikz}[european] 
		\ctikzset{voltage/distance from line=0.2} 
		\ctikzset{voltage/european label distance=2.00}  
		\draw
		(0,0) to[short] (2,0) 
		to[R,l=$R_z$] (2,-2)
		to[short] (0,-2)
		to[american current source,l=$E$] (0,0) 
		;
	\end{circuitikz}
\end{center}

Wiemy też, że w takim uproszczonym układzie będzie płynął prąd $I_1$, jak to pokazano na schemacie poniżej. 

\begin{center}\begin{circuitikz}[european] 
		\ctikzset{voltage/distance from line=0.2} 
		\ctikzset{voltage/european label distance=2.00}  
		\draw
		(0,0) to[short, i>_=$I_1$] (2,0) 
		to[R,l=$R_z$] (2,-2)
		to[short] (0,-2)
		to[american current source,l=$E$, i>_=$I_1$] (0,0) 
		;
	\end{circuitikz}
\end{center}

Wartość prądu wyznaczymy łatwo ze wzoru 
\begin{align*}
I_1=\frac{E}{R_z}
\end{align*}

Ze względu na to, że jeśli zamieniamy część układu na jego postać zastępczą, to nie zmieniają się prądy i napięcia w pozostałej części układu, wiemy, że prąd $I_1$, płynący też przez \'zródło będzie miał taką samą wartość w oryginalnym układzie:
\begin{center}\begin{circuitikz}[european] 
		\ctikzset{voltage/distance from line=0.2} 
		\ctikzset{voltage/european label distance=2.00}  
		\draw
		(0,0) to[R,l=$R_1$] (2,0) 
		to[R,l=$R_2$] (2,-2)
		to[short] (0,-2)
		to[american current source,l=$E$, i>_=$I_1$] (0,0) 
		;
	\end{circuitikz}
\end{center}

Przepływ prądu $I_1$ wywoła na rezystorach $R_1$ i $R_2$ spadki napięcia, odpowiednio 
\begin{align*}
U_1={I_1} \cdot {R_1}\\
U_2={I_1} \cdot {R_2}
\end{align*}

\begin{center}\begin{circuitikz}[european] 
		\ctikzset{voltage/distance from line=0.2} 
		\ctikzset{voltage/european label distance=2.00}  
		\draw
		(0,0) to[R,l=$R_1$, i>_=$I_1$,v^<=$U_1$] (2,0) 
		to[R,l=$R_2$, i>_=$I_1$,v^<=$U_2$] (2,-2)
		to[short] (0,-2)
		to[american current source,l=$E$, i>_=$I_1$] (0,0) 
		;
	\end{circuitikz}
\end{center}

Zatem ostatecznie napięcie na rezystorze $R_1$ wyniesie, po podstawieniu wcześniej wyznaczonych $I_1$ i $R_z$:
\begin{align*}
U_1=\frac{E}{R_1+R_2} \cdot {R_1} = E \cdot \frac{R_1}{R_1+R_2}
\end{align*}
Otrzymany wzór jest bardzo często określany jako wzór na wyjściowe napięcie dzielnika napięciowego zbudowanego z rezystorów $R_1$ i $R_2$ i warto go zapamiętać.
\end{task}